\section{Physical model and thermophysical properties}
\label{sec:numerical-framework:governing-equations}

HiSA solves the compressible conservation equations for mass, momentum and
total energy. The turbulent calculations use RANS with the eddy-viscosity and
thermal closures introduced in Section~\ref{sec:background:rans-turbulence-modelling}. The fluid is treated as a single-component or frozen-composition ideal gas, without chemical reactions. The boundary conditions and property values are chosen to represent each respective reference case and are stated in their respective section throughout this work.

\subsection{Thermophysical properties in OpenFOAM}
\label{subsec:thermophysical-properties-OF}

OpenFOAM assembles the thermophysical description from separate models for the
equation of state, caloric properties, transport properties and energy variable.
Their combination is selected in \texttt{thermophysicalProperties} and must be
registered in the installation. The \texttt{perfectGas} equation of state, \(p=\rho RT\), can be combined with constant specific heat through \texttt{hConst}, or temperature-dependent
specific heat through \texttt{janaf}. The latter uses piecewise NASA
seven-coefficient polynomials to describe the temperature dependent behaviour.
\par\medskip

Molecular viscosity \(\mu\) and thermal conductivity \(\kappa\) must reproduce
the reference laws together. Their relation is expressed by the molecular
Prandtl number, \(Pr=c_p\mu/\kappa\). In particular, the standard
\texttt{sutherland} transport model couples its viscosity law internally in OpenFOAM to a modified
Eucken conductivity relation and it is not possible to independently prescribe a constant
\(Pr\). This is a limitation of that model combination in OpenFOAM. The additional combinations described in Section~\ref{subsec:ext:thermo-physical-extension} permit separate polynomial
descriptions of \(\mu(T)\) and \(\kappa(T)\). This allows, that for the calorically perfect
benchmarks with a constant \(Pr\), the conductivity polynomial-coefficients can then be set proportional to the viscosity coefficients so that \(Pr\) can remain constant.

\subsection{Turbulent transport}

The turbulent Prandtl number \(Pr_t\) controls the modelled turbulent heat
transport and is distinct from molecular \(Pr\). With the gradient-diffusion
closure, the effective coefficients are
\(\mu_{\mathrm{eff}}=\mu+\mu_t\) and
\(\kappa_{\mathrm{eff}}=\kappa+c_p\mu_t/Pr_t\).
Unless stated otherwise, the turbulent calculations use \(Pr_t=0.9\), as
specified in Paragraph~\ref{para:turbulent-Prandtl-number}.
These effective coefficients enter both the numerical transport operations
below and the later thermal-profile assessment in
Section~\ref{subsec:comparative:zpg:wr-prandtl}.
